\pdfoutput=1

\documentclass[11pt]{article}

\usepackage[]{acl}

\usepackage{times}
\usepackage{latexsym}
\usepackage[T1]{fontenc}
\usepackage[utf8]{inputenc}
\usepackage{microtype}
\usepackage{amsmath}
\usepackage{amssymb}
\usepackage{booktabs}

\title{Bulk-Informed Cross-Modality Representation Learning\\
for scRNA-seq Batch Correction}

\author{First Author \\
  Affiliation / Address line 1 \\
  \texttt{email@domain} \\}

\begin{document}
\maketitle
\begin{abstract}
Single-cell RNA-seq (scRNA-seq) batch correction is the dominant
pre-processing step in every modern single-cell pipeline, but a recent 2025
calibration audit found that seven of eight widely used methods introduce
artefacts even under null conditions, and that evaluation metrics themselves
are unreliable. Meanwhile, bulk RNA-seq atlases (GTEx, TCGA, ENCODE) contain
orders of magnitude more donor-level coverage, cleaner technical noise, and
decades of curated biological priors---and no current integration method uses
them. We propose a cross-modality variational autoencoder with an adversarial
modality discriminator that learns a shared, technology-invariant latent space
for bulk and scRNA-seq, regularised by a \emph{gene-discordance}
supervision signal derived from matched bulk / pseudobulk pairs. The learned
representation is evaluated primarily on scIB-style integration benchmarks
(kBET, LISI, ASW, graph connectivity) and on null-condition calibration tests,
with cell-type annotation transfer and bulk deconvolution as secondary
downstream probes. The full pipeline trains in under two hours on a single
consumer GPU.
\end{abstract}

\section{Motivation}

Single-cell RNA-seq resolves tissue into individual transcriptomes but is
sparse ($\sim$80\% zeros), shallow, and subject to large technology-specific
biases: platform, chemistry version, sequencing depth, dissociation stress,
and ambient RNA. Across 15 studies PC1 correlates with detection rate
($r{=}0.54$--$0.93$), cells cluster by lab rather than by type, and the first
Alzheimer's snRNA-seq study reported 1{,}031 DEGs that collapsed to $\sim$16
after proper pseudobulk reanalysis (549-fold false-discovery inflation).
Batch correction is therefore the critical bottleneck for every downstream
scRNA-seq analysis---yet a 2025 null-condition audit shows 7 of 8 widely
deployed methods create artefacts when no batch effect exists; only Harmony
passes, and silhouette-based metrics are themselves unreliable. Foundation
models such as \citet{wen-etal-2024-cellplm} and
\citet{levine-etal-2024-cell2sentence} are trained purely on scRNA-seq and
inherit its technical biases; neither incorporates bulk at any stage.

Bulk RNA-seq is the obvious untapped prior. GTEx~v11 alone provides
$\sim$17{,}000 deeply sequenced bulk samples across 54 tissues, with cleaner
noise and no dissociation artefacts. Matched bulk/pseudobulk comparisons
expose $\sim$557 reproducibly discordant genes \citep{hu-2026-bal} whose
distributions disagree even after identical normalisation---a measurable,
gene-resolved signature of sc technical bias. No published integration
method uses bulk-derived gene-level variance, expression distributions, or
discordance signals as priors for scRNA-seq batch correction. This proposal
closes that gap.

\paragraph{Contributions.}
(i)~A dual-encoder variational architecture with an adversarial modality
classifier that aligns bulk and scRNA-seq log-CPM distributions in a shared
$d{=}64$ latent space, producing batch-corrected embeddings anchored by the
cleaner bulk modality.
(ii)~A \emph{gene-discordance} auxiliary head that learns a per-gene bulk--sc
disagreement score, supervised on matched donor pairs, which acts as an
explicit inductive bias against integrating sc batches along bias directions
bulk contradicts.
(iii)~A rigorous calibration-aware evaluation protocol: scIB integration
metrics, the 2025 null-condition random pseudo-batch test, and downstream
probes (label transfer, pseudobulk DE reproducibility vs.\ matched bulk,
deconvolution) that we expect to improve together if the representation is
honest, and disagree if any metric is gameable.
(iv)~Systematic ablations (no adversarial loss, no bulk modality, no
discordance head) that quantify which ingredient drives each metric.

\section{Problem Definition}

Let $X \in \mathbb{R}^{G \times N}$ be an scRNA-seq matrix over $G$ shared
genes and $N$ cells, with per-cell batch labels $m \in \{1,\dots,M\}^{N}$ and
optional cell-type labels $y \in \{1,\dots,K\}^{N}$. Let
$B \in \mathbb{R}^{G \times D}$ be a bulk matrix from $D$ donors (possibly
different donors from $X$). We want to learn an encoder
$E: \mathbb{R}^{G} \to \mathbb{R}^{d}$ such that, for a cell $x_i$,
$z_i = E(x_i)$ satisfies three desiderata simultaneously:
\begin{enumerate}
\itemsep=0pt
\item \textbf{Batch-invariant:} $z$ is statistically independent of $m$
  conditional on biology; formally, a classifier trained to predict $m$ from
  $z$ achieves AUC near chance on a null pseudo-batch split.
\item \textbf{Biology-preserving:} cells of the same type cluster together
  across batches; $\mathrm{ASW}(z, y)$ is high on held-out batches.
\item \textbf{Bulk-aligned:} bulk profiles embed into the \emph{same} latent
  space and nearby regions correspond to biologically matched tissues, giving
  bulk and sc a common coordinate system.
\end{enumerate}
This last requirement is what makes the problem bulk-\emph{informed} rather
than merely another sc integration model: it imports bulk's cleaner technical
noise as a prior and makes cross-modality transfer (sc $\leftrightarrow$ bulk)
a built-in capability rather than a downstream bolt-on.

\section{Method}

Three coupled modules, trained end-to-end.

\paragraph{(1) Cross-modality VAE.}
Dual encoders $E_{\mathrm{bulk}}, E_{\mathrm{sc}}: \mathbb{R}^{G} \to
\mathbb{R}^{2d}$ (MLPs, hidden $1024 \to 512$) output
$(\mu, \log\sigma^{2})$; reparameterise to $z \in \mathbb{R}^{d}$ with
$d{=}64$. A shared decoder $D: \mathbb{R}^{d} \to \mathbb{R}^{G}$
reconstructs the input. A \emph{modality discriminator}
$C_{\mathrm{mod}}: \mathbb{R}^{d} \to \{0,1\}$ (2-layer MLP) classifies $z$ as
bulk or sc; a \emph{batch discriminator} $C_{\mathrm{batch}}$ classifies $z$
among the $M$ sc batches. The encoders are trained adversarially against both
discriminators. The encoder loss is
\[
\mathcal{L}_{\mathrm{enc}}
  = \mathcal{L}_{\mathrm{recon}}
  + \beta(t)\, \mathcal{L}_{\mathrm{KL}}
  - \lambda_m(t)\, \mathcal{L}_{\mathrm{mod}}
  - \lambda_b(t)\, \mathcal{L}_{\mathrm{batch}},
\]
with $\beta, \lambda_m, \lambda_b$ annealed from $0$ to avoid posterior
collapse and early discriminator domination. When cell-type labels are
available we optionally add a scANVI-style classifier on $z$ as a
semi-supervised anchor (an ablation; the default model is unsupervised).

\paragraph{(2) Gene-discordance head.}
A small MLP $H_d : \mathbb{R}^{d} \to \mathbb{R}^{G}$ predicts, from an sc
cell's latent code $z$, a per-gene \emph{discordance score}
$\hat d_g \in \mathbb{R}$ that estimates how much that gene's value in the sc
modality disagrees with the bulk estimate of the same biology. Supervision
comes from matched donors: for each donor with paired bulk and pseudobulk, we
compute $d_{g} = |\log\mathrm{FC}(b_{g}, \mathrm{pseudobulk}(X_{g}))|$ and
regress $\hat d_g$ against $d_g$. This gives the encoder a gradient that
explicitly distinguishes \emph{genes whose sc values can be trusted} from
\emph{genes whose sc values systematically deviate from bulk}, and pushes the
latent space to downweight the latter. The head is auxiliary---it is not
required at inference.

\paragraph{(3) Bulk-prior regulariser.}
We regularise the latent distribution of sc cells toward the bulk latent
distribution at the donor / tissue level, using an MMD term
$\mathrm{MMD}^{2}\!\left(\{z^{\mathrm{sc}}_{i}\}_{y_i = t}, \{z^{\mathrm{bulk}}_{j}\}_{\mathrm{tissue}(j)=t}\right)$
summed over tissues / cell types $t$. This makes bulk's cleaner
gene-coexpression structure act as a distributional prior on the sc cells,
without requiring per-cell ground truth.

\paragraph{Training schedule.}
\emph{(a)~Pretraining:} encoders, decoder, and both discriminators on
unpaired bulk (GTEx, $\sim$800 donors) and sc (HCA, $\sim$390k cells) with
only the VAE + adversarial losses, to shape an aligned latent space.
\emph{(b)~Fine-tuning:} discordance head and MMD regulariser added, trained
jointly on paired matched donors. No pretrained foundation models are used;
all parameters ($\sim$6M total) are trained from scratch.

\paragraph{Why each component is needed.}
A naive VAE on the union of bulk and sc cleanly separates the two into
disjoint latent clusters (the reconstruction objective alone gives no
alignment pressure). The discordance head supplies an explicit gradient
distinguishing biological batch differences from technology-specific ones
(the 2025 audit's central failure mode). The MMD term imports the bulk
prior at the \emph{distributional} level, which the binary discriminator
cannot. The three components target three distinct failure modes.

\section{Datasets and Evaluation Protocol}

\subsection{Datasets}

\textbf{Pretraining (unpaired).}
\emph{Bulk:} GTEx v11 whole-blood subset, $\sim$803 donors, filtered to
$\sim$16k genes expressed in $\geq$10\% of donors at CPM~$\geq$~1.
\emph{Single-cell:} HCA blood scRNA-seq, $\sim$390k cells, down-sampled to a
balanced 50k-cell pretraining set stratified across donors and cell types.

\textbf{Batch-correction benchmark (primary evaluation).}
A scIB-style multi-batch cohort constructed from HCA blood and published
10x/Smart-seq2 pairs. Each cell has batch, donor, and cell-type labels, so
both batch-mixing and biology-preservation can be scored.

\textbf{Held-out matched bulk/sc cohort (discordance ground truth).}
The BAL paired benchmark of \citet{hu-2026-bal} (30 donors with matched bulk
and sc) supplies the 557 published discordant genes used as validation for
the gene-discordance head.

\textbf{Downstream probe datasets.}
\emph{(i)~Cell-type transfer:} pancreatic islet scRNA-seq (GSE84133, Baron et
al.) across 14 donors, label-transferred from a held-out donor.
\emph{(ii)~Pseudobulk DE vs.\ real bulk:} sc $\to$ pseudobulk DE on a held-out
cohort, compared to matched real bulk DE via rank concordance.
\emph{(iii)~Bulk deconvolution:} GSE50244 pancreas bulk (89 donors with
HbA1c) deconvolved with an NNLS solver applied on top of signatures derived
from the corrected sc embeddings, reported as a sanity check that the latent
space remains quantitatively linked to expression magnitudes.

\subsection{Evaluation Protocol}

\paragraph{Splits.} Donors and batches are split 70/15/15 train/val/test. All
cells from a test donor or a test batch are \emph{never} seen during training
(donor- \emph{and} batch-level leakage control).

\paragraph{Primary metrics (mean $\pm$ std, 5 seeds).}
Batch mixing: kBET, iLISI, graph connectivity, batch ASW. Biology
conservation: cLISI, cell-type ASW, Leiden NMI/ARI, isolated-label F1. We
also report the scIB composite (since it alone masks overcorrection) and the
2025 null-condition calibration test (fraction of spurious clusters under a
random pseudo-batch split---an honest integrator creates none).

\paragraph{Secondary / downstream metrics.}
Cell-type label transfer accuracy and macro-F1 from held-out donors;
pseudobulk-vs-real-bulk DE rank concordance (Spearman on $|\log\mathrm{FC}|$)
and discordant-gene recall on the BAL cohort; deconvolution RMSE/MAE/CCC on
GSE50244 as a quantitative-fidelity probe (not the primary metric).

\paragraph{Baselines (fair comparison).}
All baselines consume the \emph{same} preprocessed expression matrix, share
the same train/val/test splits, and are evaluated with identical metrics on
identical test sets: Harmony (the 2025 audit's only well-calibrated method),
Seurat v5 RPCA, scVI, scANVI, Scanorama, and ComBat-seq. Ablations isolate
each contribution: (a)~no adversarial modality loss, (b)~no bulk modality
(sc-only encoder), (c)~no discordance head, (d)~no MMD regulariser.

\section{Estimated Computing Resources}

\textbf{Hardware.} One consumer GPU (NVIDIA RTX 3070/4070 class, 8--12\,GB
VRAM) or a free Google Colab T4. \textbf{Memory.} The full HCA blood
sub-sample (50k cells $\times$ 16k genes at fp32) is $\sim$3\,GB.
\textbf{Runtime.}
(i)~Data preprocessing: $<$10\,min on CPU.
(ii)~Cross-modality VAE pretraining: $\sim$60\,min (100 epochs, batch 256).
(iii)~Joint fine-tuning with discordance and MMD heads: $\sim$30\,min.
(iv)~Full 5-seed sweep across all baselines, metrics, and ablations:
$<$10 GPU-hours end-to-end.
\textbf{Storage.} Raw datasets $\sim$10\,GB; processed artefacts $\sim$1.5\,GB.

The project uses \emph{no} pretrained foundation models, so compute does not
scale with any external checkpoint. All code is PyTorch, scikit-learn,
scanpy, and scib, already installed and exercised in the existing project
repository. The end-to-end pipeline is comfortably executable within the
remaining course window with headroom for additional ablations.

\paragraph{Timeline.}
\emph{Weeks 1--2:} assemble the scIB cohort, reproduce Harmony/scVI/scANVI
baselines. \emph{Weeks 3--4:} cross-modality VAE + adversarial trainer;
verify null-condition calibration. \emph{Weeks 5--6:} discordance head + MMD
regulariser + ablations. \emph{Week 7:} downstream probes, paired bootstrap
tests, write-up. A one-week buffer remains before the deadline.

\paragraph{Risks and fallbacks.}
If adversarial training is unstable, swap the discriminator for pure MMD
alignment. If the discordance head overfits the small matched cohort, freeze
it as a fixed gene weighting. Both fallbacks preserve the core cross-modality
integration contribution.

\section{Related Work}

\textbf{Single-cell foundation models.}
\citet{wen-etal-2024-cellplm} pretrain a cell language model on scRNA-seq
with strong imputation and perturbation-prediction performance. The model is
scRNA-seq only: bulk RNA-seq, with its cleaner noise and orders-of-magnitude
larger donor coverage, is not used at any stage. The same is true of every
foundation model in the 2025 survey. Our work is complementary: bulk is a
first-class modality during pretraining, and the shared latent space is
explicitly aligned by an adversarial loss rather than left to emerge.

\textbf{LLMs for single-cell.}
\citet{levine-etal-2024-cell2sentence} tokenise scRNA-seq profiles as ranked
gene sequences and apply an LLM. The gene-rank tokenisation discards absolute
expression magnitudes that matter for batch correction (where detection rate
itself is the dominant axis of variation). Our method keeps a continuous
latent representation, so magnitude information and gradient-based integration
losses remain available.

\textbf{Cross-modal representation learning.}
\citet{xu-etal-2023-mmt} survey alignment, fusion, and translation strategies
for multimodal transformers, primarily in vision-language. We adopt their
fusion-via-shared-latent paradigm but replace cross-attention alignment
(appropriate for token-aligned modalities) with an adversarial modality
classifier and an MMD distributional prior, because bulk and scRNA-seq have
no natural token correspondence and differ primarily in their
\emph{distribution}, not in their feature layout.

\paragraph{Justified novelty.}
Relative to the three works above, our contribution is the first to combine
the following in a single end-to-end batch-correction model: (i)~bulk and
scRNA-seq as first-class modalities in a shared-latent VAE, beyond the
sc-only pretraining of CellPLM; (ii)~an explicitly supervised gene-discordance
head anchored on measurable bulk/pseudobulk disagreement, addressing the
calibration failure mode ignored by both CellPLM and Cell2Sentence; and
(iii)~a distributional MMD prior that imports bulk's cleaner noise structure
into the sc integration objective, a use of the shared-latent paradigm of
\citet{xu-etal-2023-mmt} that their vision-language setting does not need.
Each ingredient has precedent in isolation, but the combination, targeted at
the calibration crisis of scRNA-seq batch correction and anchored on the
discordance signal that matched bulk/pseudobulk pairs make measurable, has
not been reported in any of the allowed venues.

\bibliography{anthology,custom}
\bibliographystyle{acl_natbib}

\end{document}
